\chapter*{Abstrak}
\vspace*{0.2cm}

{
\setlength{\parindent}{0pt}
\noindent
Nama: \penulis\par
Program Studi: \program\par
Judul: \judul\par
\medskip

\noindent
Pertumbuhan data teks dari platform digital, seperti ulasan produk dan media
sosial, mendorong kebutuhan akan metode otomatis yang mampu menganalisis
sentimen dan mengidentifikasi topik secara bersamaan. Pendekatan analisis
sentimen konvensional umumnya mengasumsikan satu sentimen untuk setiap
dokumen sehingga belum mampu merepresentasikan ulasan yang membahas lebih
dari satu topik. Penelitian ini mengembangkan model \textit{Joint
Document-Level Sentiment--Topic} berbasis \textit{Deep Embedded Clustering}
(JDST-DEC) yang mengintegrasikan analisis sentimen dan pemodelan topik
secara \textit{end-to-end} dalam kerangka \textit{Multi-Task Learning}
(MTL), menggunakan satu \textit{shared encoder} berbasis BERT dan DEC
sebagai \textit{topic head} untuk membentuk representasi topik dari
representasi laten yang dihasilkan encoder. Pelatihan dilakukan dalam dua
tahap, yaitu pra-pelatihan \textit{autoencoder} untuk membentuk representasi
laten, kemudian optimisasi gabungan menggunakan \textit{sentiment loss}
berbasis \textit{cross-entropy} dan \textit{clustering loss} berbasis KL
\textit{divergence}. Evaluasi dilakukan pada dataset IMDB dan Yelp dengan
variasi jumlah topik $K=30$, $50$, dan $80$. Pada tugas klasifikasi
sentimen, JDST-DEC memperoleh F1-\textit{score} terbaik sebesar 0{,}8986
pada IMDB dan 0{,}9020 pada Yelp ($K=50$), lebih tinggi dibandingkan
NN-BERT (0{,}8612 dan 0{,}8715). Pada tugas pemodelan topik, JDST-DEC
menghasilkan kualitas topik yang lebih baik dibandingkan MTL-GUI pada
beberapa konfigurasi, namun masih berada di bawah BERTopic-KMeans
berdasarkan metrik \textit{Topic Coherence} (TC-NPMI) dan \textit{Topic
Diversity}. Hasil ini menunjukkan bahwa \textit{Deep Embedded Clustering}
berhasil diintegrasikan ke dalam kerangka JDST dan mampu meningkatkan
performa klasifikasi sentimen dibandingkan model \textit{single-task},
meskipun kualitas representasi topik masih berada di bawah metode
pemodelan topik khusus.
\medskip

\noindent
\textbf{Kata kunci}: \textit{Deep Embedded Clustering},
\textit{Joint Document-Level Sentiment--Topic Modeling},
\textit{Multi-Task Learning}, \textit{BERT}.
}

\newpage
